\documentclass[12pt]{article}

\usepackage{graphicx}
\usepackage{url}
\usepackage{hyperref}
\usepackage{verbatim}
\usepackage[title,titletoc,toc]{appendix}
\usepackage{wasysym}

\renewcommand{\topfraction}{1.0}
\renewcommand{\bottomfraction}{1.0}
\renewcommand{\textfraction}{0.0}
\renewcommand{\floatpagefraction}{1.0}
\renewcommand{\dbltopfraction}{1.0}


\begin{document}

\begin{titlepage}

\begin{center}
{\large STAT 5810/6810 Introduction to Statistical Computing} \\[4cm]

{\LARGE \bf Reproducible Research (RR)} \\[1cm]
by \\[0.5cm]
{\bf J\"urgen Symanzik} \\[2.5cm]
{\bf Date:} \today \\[2cm]

UTAH STATE UNIVERSITY \\[0.5cm]
Logan, UT \\[0.5cm]
Fall 2012 \\[0.5cm]
\end{center}

\thispagestyle{empty}
\vfill
\end{titlepage}


\newpage 


\pagenumbering{roman}

\tableofcontents


\newpage

\pagenumbering{arabic}


\section{Basics about \LaTeX\ via RStudio}

First, notice that you can open a \LaTeX\ .tex file from within RStudio.
This file is recognized as a .tex file --- and you can translate it 
via the Compile PDF option from within RStudio!

When working with .tex in RStudio, you have options to introduce
some \LaTeX\ formatting via highlighting and mouse clicke. Let's
try this for the words \textbf{bold}, \emph{italic}, and \texttt{typewriter}.


\subsection*{Now, we want a new subsection}

Notice what is being produced here!



And what about a bullet list hereafter?
\begin{itemize}
  \item xxx
\end{itemize}



\section{Using Sweave and knitr in RStudio}

Detailed information can be found at
\url{http://www.rstudio.com/ide/docs/authoring/overview}.

Basically, you need to start at \textbf{File $->$ New $->$ R Sweave}.
Let's do so.


If you want to use an existing .tex file, you must change the file
extension to .Rnw and add the line
\begin{verbatim}
\SweaveOpts{concordance=TRUE}
\end{verbatim}
after the 
\begin{verbatim}
\begin{document}
\end{verbatim}
line.
 

\end{document}

